\subsection{可复用表达与“避雷”边界}

可复用的是信息结构，不是整句套用。摘要可用“针对对象和任务—建立核心模型—
给出关键结果—报告检验”的四步；模型段可用“关系依据—变量定义—方程/约束—
式后解释”；结果段可用“读数—比较—机制—决策”；评价段可用“机制带来的能力—
假设造成的边界—可验证改进”。“首先、其次、最后”“通过软件求解可得”“结果
较为理想”“模型具有较强普适性”等句子若没有对象、数值、条件和证据，应删除
或改写。

表达复用不等于隐藏不确定性。OCR 无法确认的字词标为未确认；跨论文统计只报告
检测定义和样本数；本地缺少的 2020 赛题原文不由评述补写。论文中的个别错误也
不应遮蔽其优秀部分：分析重点是任务闭合、关键推导、结果解释和验证等产生高
判别力的做法，同时把错误作为质量门输入。

